% !TEX root = ../main.tex
\corrigesChapitre{9}{Force des acides et des bases}
\corrigesSouples

\begin{blocLivre}{cor-9-1}
\begin{corrige}{1}{Le produit ionique de l'eau}
\demarche{Le produit ionique relie les deux concentrations : connaître l'une, c'est connaître l'autre.}
\begin{questions}
\item \ce{2 H2O(l) <=> H3O^+(aq) + HO^-(aq)} et $K_\mathrm{e}=\dfrac{[\ce{H3O+}]_\text{éq}\,[\ce{HO-}]_\text{éq}}{(\cst)^2}$.
\item $[\ce{H3O+}]=10^{-2{,}0}=\qty{1,0e-2}{mol.L^{-1}}$ et $[\ce{HO-}]=\dfrac{K_\mathrm{e}}{10^{-2{,}0}}=\res{\qty{1,0e-12}{mol.L^{-1}}}$.
\item $[\ce{H3O+}]=\qty{1,0e-11}{mol.L^{-1}}$, donc $[\ce{HO-}]=\res{\qty{1,0e-3}{mol.L^{-1}}}$.
\item Dans une eau neutre, $[\ce{H3O+}]=[\ce{HO-}]$, donc $[\ce{H3O+}]^2=K_\mathrm{e}(\cst)^2$ et
$\pH=\frac12\,\mathrm{p}K_\mathrm{e}$. À 60~°C, $\mathrm{p}K_\mathrm{e}=-\log(\num{9.6e-14})=13{,}0$ et $\pH=\res{6{,}51}$.
L'eau reste \res{\text{neutre}} : les deux concentrations sont égales, même si le pH est inférieur à 7.
\end{questions}
\end{corrige}
\end{blocLivre}
\begin{blocLivre}{cor-9-2}
\begin{corrige}{2}{Acide fort ou acide faible ?}
\demarche{Comparer la concentration en ions oxonium, donnée par le pH, à la concentration apportée.}
\begin{questions}
\item Un acide fort réagit totalement avec l'eau ; un acide faible réagit de façon limitée : il subsiste des molécules
d'acide dans l'état final.
\item \ce{HCl + H2O -> Cl- + H3O+} (flèche simple) ; \ce{CH3COOH + H2O <=> CH3COO- + H3O+} (double flèche).
\item HCl : $\tau=\dfrac{10^{-2{,}0}}{\num{1.0e-2}}=\res{1}$ ; acide éthanoïque :
$\tau=\dfrac{10^{-3{,}4}}{\num{1.0e-2}}=\res{0{,}040}$. Le premier est fort, le second faible.
\item La solution \res{\text{d'acide éthanoïque}} : 96~\% des molécules n'ont pas réagi, alors que la solution d'acide
chlorhydrique n'en contient plus.
\end{questions}
\end{corrige}
\end{blocLivre}
\begin{blocLivre}{cor-9-3}
\begin{corrige}{3}{Écrire des constantes d'acidité}
\demarche{Le pKa diminue quand la force de l'acide augmente ; les acides du tableau de l'essentiel servent de repères.}
\begin{questions}
\item Pour chaque couple, $\ce{AH + H2O <=> A- + H3O+}$ et $K_\mathrm{A}=\dfrac{[\ce{A-}]_\text{éq}[\ce{H3O+}]_\text{éq}}{[\ce{AH}]_\text{éq}\,\cst}$,
avec (AH ; \ce{A-}) valant (\ce{HCOOH} ; \ce{HCOO-}), (\ce{NH4+} ; \ce{NH3}) ou (\ce{HClO} ; \ce{ClO-}).
\item \ce{HCOOH} : $\pKa=3{,}75$, $K_\mathrm{A}=10^{-3{,}75}=\res{\num{1.8e-4}}$ ; \ce{HClO} : 7,5,
$K_\mathrm{A}=\res{\num{3.2e-8}}$ ; \ce{NH4+} : 9,2, $K_\mathrm{A}=\res{\num{6.3e-10}}$.
\item Du plus fort au plus faible : $\res{\ce{HCOOH}>\ce{HClO}>\ce{NH4+}}$.
\end{questions}
\end{corrige}
\end{blocLivre}
\begin{blocLivre}{cor-9-4}
\begin{corrige}{4}{Lire un diagramme de prédominance}
\demarche{Comparer le pH de chaque boisson au pKa : en dessous, la forme acide prédomine.}
\begin{questions}
\item Voir le diagramme complété ci-dessous.
\item Soda ($\pH=2{,}5<4{,}2$) : \res{\ce{C6H5COOH}} ; jus de tomate ($\pH=4{,}2=\pKa$) : les deux formes à égalité ;
eau pétillante ($\pH=5{,}5>4{,}2$) : \res{\ce{C6H5COO-}}.
\item $\dfrac{[\ce{C6H5COOH}]}{[\ce{C6H5COO-}]}=10^{\,\pKa-\pH}=10^{1{,}7}=\res{50}$ : la forme basique est
ultraminoritaire (2~\% environ).
\item Dans le \res{\text{jus de tomate}}, dont le pH est égal au $\pKa$.
\end{questions}
\schema[corrige]{09-ex04}
\end{corrige}
\end{blocLivre}
\begin{blocLivre}{cor-9-5}
\begin{corrige}{5}{Du pH à la concentration}
\demarche{Acide fort : $\pH=-\log(c/\cst)$ ; base forte : $\pH=14+\log(c/\cst)$ ; les relations s'inversent avec les puissances de 10.}
\begin{questions}
\item $\pH=-\log(\num{2.0e-3})=\res{2{,}70}$.
\item $\pH=14{,}0+\log(\num{1.0e-3})=\res{11{,}0}$.
\item $[\ce{HO-}]=10^{\,\pH-14}=10^{-3{,}6}=\res{\qty{2,5e-4}{mol.L^{-1}}}$.
\item Le pH augmente d'une unité si la concentration est divisée par 10, donc si le volume est multiplié par 10 :
il faut atteindre 100~mL, soit ajouter \res{\qty{90}{mL}} d'eau.
\end{questions}
\end{corrige}
\end{blocLivre}
\begin{blocLivre}{cor-9-6}
\begin{corrige}{6}{Comparer la force de deux acides}
\demarche{Un pKa plus petit signifie un acide plus fort : à même concentration, pH plus faible et taux plus grand.}
\begin{questions}
\item L'acide méthanoïque a le $\pKa$ le plus petit : il est plus fort, son taux d'avancement final est plus grand et
sa solution a le \res{\text{pH le plus faible}}.
\item À l'équilibre, $[\ce{HCOO-}]=[\ce{H3O+}]=h$ et $[\ce{HCOOH}]=c-h$, d'où $K_\mathrm{A}=\frac{h^2}{c-h}$ et
$h^2+K_\mathrm{A}h-K_\mathrm{A}c=0$. La racine positive vaut
\[h=\frac{-K_\mathrm{A}+\sqrt{K_\mathrm{A}^2+4K_\mathrm{A}c}}{2}=\res{\qty{1,25e-3}{mol.L^{-1}}},\qquad\pH=\res{2{,}90}.\]
\item $\tau_{\text{méth}}=\frac{\num{1.25e-3}}{\num{1.0e-2}}=\res{12{,}5~\%}$ ;
$\tau_{\text{éth}}=\frac{10^{-3{,}39}}{\num{1.0e-2}}=\res{4{,}1~\%}$. L'acide au $\pKa$ le plus petit a bien le taux le
plus grand.
\item La dilution favorise la dissociation (chapitre~7) : les deux taux d'avancement finals \res{\text{augmentent}}.
\end{questions}
\end{corrige}
\end{blocLivre}
\begin{blocLivre}{cor-9-7}
\begin{corrige}{7}{Exploiter un diagramme de distribution}
\demarche{À pH faible, la forme acide domine : la courbe décroissante est celle de la forme acide.}
\begin{questions}
\item À pH faible, l'acide prédomine : la courbe \res{\text{décroissante}} correspond à \ce{CH3COOH}, la croissante à
\ce{CH3COO-}.
\item Les courbes se croisent (50~\%/50~\%) à $\pH=\pKa\approx\res{4{,}8}$.
\item À pH 3,8 : \res{90~\%} de \ce{CH3COOH} et \res{10~\%} de \ce{CH3COO-}, d'où un rapport voisin de $0{,}11$.
\item $\dfrac{[\ce{CH3COO-}]}{[\ce{CH3COOH}]}=10^{\,\pH-\pKa}=10^{-0{,}96}=\res{0{,}11}$ : la lecture est confirmée.
\item À pH 2,8, $\alpha(\ce{CH3COO-})=1~\%$ : l'acide éthanoïque du vinaigre est presque exclusivement sous
\res{\text{forme acide }\ce{CH3COOH}}.
\end{questions}
\schema[corrige]{09-ex07}
\end{corrige}
\end{blocLivre}
\begin{blocLivre}{cor-9-8}
\begin{corrige}{8}{Le pH d'une solution d'acide méthanoïque}
\demarche{Le tableau d'avancement donne les concentrations à l'équilibre en fonction de h ; Ka fournit l'équation.}
\begin{questions}
\item Pour un volume $V$ : $n(\ce{HCOOH})=cV-x$, $n(\ce{HCOO-})=n(\ce{H3O+})=x$ ; à l'équilibre, $h=x_\mathrm{f}/V$.
\item $K_\mathrm{A}=\dfrac{h^2}{c-h}$ donne $h^2+K_\mathrm{A}h-K_\mathrm{A}c=0$, de racine positive
\[h=\frac{-\num{1.8e-4}+\sqrt{(\num{1.8e-4})^2+4\times\num{1.8e-4}\times\num{5.0e-3}}}{2}=\res{\qty{8,6e-4}{mol.L^{-1}}},\]
d'où $\pH=-\log(\num{8.63e-4})=\res{3{,}06}$.
\item $\tau=\dfrac{h}{c}=\dfrac{\num{8.63e-4}}{\num{5.0e-3}}=\res{17{,}3~\%}$.
\item $h\approx\sqrt{K_\mathrm{A}c}=\qty{9,5e-4}{mol.L^{-1}}$, soit $\pH\approx\res{3{,}02}$. L'écart (0,04 unité) n'est
pas négligeable : avec un taux de dissociation de 17~\%, l'acide n'est pas « peu dissocié » et l'approximation
surestime $h$.
\end{questions}
\end{corrige}
\end{blocLivre}
\begin{blocLivre}{cor-9-9}
\begin{corrige}{9}{Le bleu de bromothymol}
\demarche{La couleur observée est celle de la forme prédominante ; autour du pKa, les couleurs se mélangent.}
\begin{questions}
\item Les deux formes du couple $\mathrm{HIn}/\mathrm{In}^-$ ont des couleurs différentes et leur proportion dépend du
pH : la couleur de la solution change donc avec le pH.
\item Suc gastrique ($\pH=2<6{,}0$) : \res{\text{jaune}} ; eau de mer ($8{,}2>7{,}6$) : \res{\text{bleu}} ; eau pure
($\pH=7{,}0$, dans la zone de virage) : \res{\text{vert}}, mélange des deux couleurs.
\item $\dfrac{[\ce{In-}]}{[\ce{HIn}]}=10^{\,7{,}4-7{,}1}=10^{0{,}3}=\res{2{,}0}$.
\item La couleur d'une forme ne s'impose que si elle est environ dix fois plus concentrée que l'autre, c'est-à-dire
pour $\pH<\pKa-1$ ou $\pH>\pKa+1$ : la zone de virage s'étend sur \res{\pKa\pm1\text{ environ}}.
\item L'indicateur est lui-même un couple acide-base : en quantité notable, il \res{\text{modifierait le pH}} de la
solution étudiée.
\end{questions}
\end{corrige}
\end{blocLivre}
\begin{blocLivre}{cor-9-10}
\begin{corrige}{10}{Une solution tampon}
\demarche{Tant que les quantités des deux formes restent comparables, le pH reste proche du pKa.}
\begin{questions}
\item $[\ce{CH3COO-}]_\text{éq}=[\ce{CH3COOH}]_\text{éq}$, donc $\pH=\pKa+\log1=\pKa=\res{4{,}76}$.
\item Le chlorure d'hydrogène transforme \ce{CH3COO-} en \ce{CH3COOH} :
$n(\ce{CH3COO-})=\qty{0,099}{mol}$ et $n(\ce{CH3COOH})=\qty{0,101}{mol}$, d'où
\[\pH=4{,}76+\log\frac{0{,}099}{0{,}101}=\res{4{,}75}.\]
\item Dans l'eau pure, l'acide fort impose $[\ce{H3O+}]=\qty{1,0e-3}{mol.L^{-1}}$ : $\pH=\res{3{,}00}$.
\item Le pH de S ne varie que de 0,01, contre 4 unités dans l'eau (de 7,0 à 3,0) : S est une
\res{\text{solution tampon}}.
\item Le maintien du pH du sang à 7,4 : une variation de quelques dixièmes d'unité est dangereuse.
\end{questions}
\end{corrige}
\end{blocLivre}
\begin{blocLivre}{cor-9-11}
\begin{corrige}{11}{Tracer un diagramme de distribution}
\demarche{La relation $\pH=\pKa+\log\left([\ce{A-}]/[\ce{AH}]\right)$ donne le rapport des deux formes, donc leur proportion.}
\begin{questions}
\item $\dfrac{[\ce{A-}]}{[\ce{AH}]}=10^{\,\pH-\pKa}$, donc
$\alpha=\dfrac{[\ce{A-}]}{[\ce{A-}]+[\ce{AH}]}=\dfrac{10^{\,\pH-\pKa}}{1+10^{\,\pH-\pKa}}=\res{\dfrac{1}{1+10^{\,\pKa-\pH}}}$.
\item À $\pKa-1$, la forme basique représente \res{9{,}1~\%} ; à $\pH=\pKa$, \res{50~\%} ; à $\pKa+1$,
\res{90{,}9~\%}.
\item À $\pKa+1$, la forme acide représente $1-0{,}909=\res{9{,}1~\%}$. Les deux courbes se déduisent l'une de
l'autre par symétrie par rapport à la droite verticale $\pH=\pKa$ : $\alpha(\pKa+x)=1-\alpha(\pKa-x)$.
\item Avec \code{pKa = 6.1} et \code{print(round(1 - alpha(7.4), 3))}, le programme affiche \res{\code{0.048}} : à
pH 7,4, 4,8~\% du couple est sous forme \ce{CO2,H2O}.
\end{questions}
\schema[corrige]{rev-09-ex11-cor}
\legende{Les deux formes totalisent 100 \% ; elles sont en proportions égales pour pH = pKa.}
\end{corrige}
\end{blocLivre}
\begin{blocLivre}{cor-9-12}
\begin{corrige}{12}{Le pH du sang}
\demarche{La relation pH = pKa + log(b/a) relie directement le pH au rapport des deux formes du tampon.}
\begin{questions}
\item $\res{\ce{CO2,H2O + H2O <=> HCO3^- + H3O+}}$.
\item $\pH=7{,}4>\pKa=6{,}1$ : les ions \res{\ce{HCO3-}\text{ prédominent}}.
\item $\dfrac{[\ce{HCO3-}]}{[\ce{CO2,H2O}]}=10^{\,7{,}4-6{,}1}=10^{1{,}3}=\res{20}$, d'où
$[\ce{CO2,H2O}]=\dfrac{24}{20}=\res{\qty{1,2}{mmol.L^{-1}}}$.
\item L'effort produit du \ce{CO2} (et de l'acide lactique) : $[\ce{CO2,H2O}]$ augmente, donc le rapport
$[\ce{HCO3-}]/[\ce{CO2,H2O}]$ diminue et le pH baisse. La ventilation accélérée élimine du \ce{CO2} et ramène le pH
vers 7,4.
\item Si $[\ce{CO2,H2O}]$ est divisée par 2 à $[\ce{HCO3-}]$ constante, le rapport double et
$\Delta\pH=\log2=\res{0{,}30}$ : le pH monte vers 7,7. C'est l'alcalose respiratoire du document~2, provoquée par
l'hyperventilation.
\end{questions}
\schema[corrige]{rev-09-ex12-cor}
\legende{La position du pH étudié permet de justifier la forme prédominante.}
\end{corrige}
\end{blocLivre}
\begin{blocLivre}{cor-9-13}
\begin{corrige}{13}{Le degré Dornic d'un lait}
\demarche{Convertir les degrés Dornic en concentration molaire, puis appliquer la méthode de l'acide faible.}
\begin{questions}
\item $16<18$ : le lait est \res{\text{frais}}. Il contient $16\times0{,}10=\qty{1,6}{g.L^{-1}}$ d'acide lactique, soit
$c=\dfrac{1{,}6}{90{,}1}=\res{\qty{1,8e-2}{mol.L^{-1}}}$.
\item $h^2+K_\mathrm{A}h-K_\mathrm{A}c=0$ avec $K_\mathrm{A}=\num{1.4e-4}$ : $h=\qty{1,5e-3}{mol.L^{-1}}$, d'où
$\pH=\res{2{,}82}$ et $\tau=\dfrac{h}{c}=\res{8{,}5~\%}$.
\item $\pH=6{,}7>\pKa=3{,}86$ : l'acide lactique du lait est surtout sous forme d'\res{\text{ions lactate}}. Le lait
n'est pas une simple solution d'acide lactique : ses protéines, phosphates et hydrogénocarbonates tamponnent le milieu
et neutralisent une grande partie de l'acide formé.
\item $\dfrac{[\text{lactate}]}{[\text{acide}]}=10^{\,4{,}5-3{,}86}=10^{0{,}64}=\res{4{,}4}$ : la forme basique domine,
mais la forme acide représente encore 19~\% du couple, ce qui reste cohérent avec le goût acide du yaourt.
\end{questions}
\end{corrige}
\end{blocLivre}
\begin{blocLivre}{cor-9-14}
\begin{corrige}{14}{Un conservateur et son efficacité}
\demarche{La fraction active est la proportion de la forme acide, $1/\left(1+10^{\,\pH-\pKa}\right)$.}
\begin{questions}
\item $c_{\max}=\dfrac{0{,}150}{122{,}1}=\res{\qty{1,23e-3}{mol.L^{-1}}}$. La solubilité (2,9~g$\cdot$L$^{-1}$, soit
près de vingt fois 150~mg$\cdot$L$^{-1}$) permet largement de l'atteindre.
\item À pH 3,2 : $f=\dfrac{1}{1+10^{\,3{,}2-4{,}2}}=\dfrac{1}{1+0{,}1}=\res{0{,}91}$ : 91~\% du conservateur est actif.
\item À pH 5,5 : $f=\dfrac{1}{1+10^{1{,}3}}=\res{0{,}048}$, soit 4,8~\%. Le conservateur est environ
\res{\text{19 fois moins efficace}} dans le thé glacé que dans le soda.
\item À pH 6,0, $f=\dfrac{1}{1+10^{1{,}8}}=0{,}016$ : moins de 2~\% du conservateur serait actif, même à la dose
maximale. Il faut \res{\text{déconseiller l'acide benzoïque}} pour cette boisson : abaisser le pH ou choisir un
conservateur efficace à pH plus élevé.
\end{questions}
\end{corrige}
\end{blocLivre}
\begin{blocLivre}{cor-9-15}
\begin{corrige}{15}{Désinfecter l'eau d'une piscine}
\demarche{La fraction d'acide hypochloreux se lit sur la courbe, puis se calcule avec $1/\left(1+10^{\,\pH-\pKa}\right)$.}
\begin{questions}
\item \ce{HClO + H2O <=> ClO- + H3O+} et $K_\mathrm{A}=\dfrac{[\ce{ClO-}]_\text{éq}[\ce{H3O+}]_\text{éq}}{[\ce{HClO}]_\text{éq}\,\cst}$.
\item Lecture : environ \res{67~\%} à pH 7,2 et \res{17~\%} à pH 8,2.
\item $f=\dfrac{1}{1+10^{\,\pH-7{,}5}}$ : à 7,2, $f=\dfrac{1}{1+10^{-0{,}3}}=\res{0{,}67}$ ; à 8,2,
$f=\dfrac{1}{1+10^{0{,}7}}=\res{0{,}17}$.
\item À pH 8,2, la forme active ne représente plus que 17~\% du chlore : la désinfection devient insuffisante et les
algues se développent. Il faut \res{\text{abaisser le pH}} vers 7,2--7,4 avec un correcteur « pH moins » : la fraction
d'acide hypochloreux remonte vers 67~\% sans ajout de chlore.
\item À pH 7,2--7,4, la fraction active reste élevée (56 à 67~\%) tout en gardant un pH proche de celui de la peau et
des yeux : c'est le compromis efficacité-confort du document~2.
\end{questions}
\schema[corrige]{09-ex15}
\end{corrige}
\end{blocLivre}
